%% CAPITULO 1 -- INTRODUCAO (esqueleto da dissertação; base: pré-projeto v14)

\section{Introdução}

Órgãos públicos produzem e consultam diariamente leis, decretos, portarias,
manuais e guias cuja interpretação depende da estrutura do próprio documento:
uma regra pode estar separada de sua exceção, e um artigo, de seus parágrafos.
Sistemas de geração aumentada por recuperação (RAG) prometem respostas
fundamentadas nesses documentos ao combinar a memória paramétrica de um modelo
gerador com um índice de trechos recuperados, arquitetura que, segundo
\citeonline{lewis2020rag}, produz linguagem mais específica, diversa e factual
do que a de um modelo apenas paramétrico. A qualidade dessas respostas, porém,
depende de como os documentos são segmentados e contextualizados antes da
recuperação.

Ao implantar um sistema RAG sobre documentos públicos, o órgão precisa escolher
entre estratégias de contextualização documental --- \textit{chunking} fixo,
recursivo, semântico, recuperação \textit{parent-child} sensível à estrutura e
uso de metadados --- sem um instrumento que compare, ao mesmo tempo, qualidade
da resposta, custo, rastreabilidade da fonte e risco institucional. Esta
dissertação parte, por isso, da seguinte pergunta: como propor e aplicar um
modelo multidimensional para avaliar estratégias de contextualização documental
em sistemas RAG aplicados a documentos institucionais públicos? O objetivo geral
é propor e aplicar esse modelo, cuja contribuição principal é o Índice de
Rastreabilidade Documental (IRD), validado como instrumento de medida e
acompanhado de uma matriz de decisão para o comitê de governança de dados e de
inteligência artificial do órgão.
